% ==============================================================================
% SECTION 1 — INTRODUCTION
% ==============================================================================
\section{Introduction}

When \LLM-powered agents interact repeatedly in competitive settings, do
they cooperate or defect? The question is not merely theoretical: autonomous
\LLM agents are already deployed to negotiate contracts, allocate computational
resources, and bid in markets \cite{Wang2024_survey}. In each of these settings,
the long-run social welfare of the system hinges on whether evolutionary pressure
selects for cooperative or aggressive behaviour.

\citet{Willis2025_llm_ipd} provided the first systematic treatment
of this question, using the Iterated Prisoner's Dilemma (\IPD) as a formal
testbed. Rather than prompting \LLM{}s to output individual actions---an
approach that prior work found unreliable \cite{Fan2024_rational}---they prompt
models to generate \emph{complete strategies} in natural language, which are
subsequently implemented as Python algorithms. By simulating populations of
agents through a Moran evolutionary process, they characterised the equilibrium
tendencies of ChatGPT-4o and Claude~3.5~Sonnet. Their central finding is a
persistent cooperative bias: in balanced populations and under most prompting
conditions, cooperative strategies dominate, with aggressive equilibria arising
at substantially below prior probability.

Two gaps limit the reach of that benchmark. First, the \LLM landscape has
shifted substantially since that study. A new generation of frontier models,
released throughout 2025--2026, incorporates extended context, revised alignment
training, and improved reasoning capabilities. Whether their cooperative
tendencies persist, deepen, or erode under evolutionary pressure is unknown.
Second, the original study covered only two models from two providers, leaving
it unclear whether cooperative bias is a universal property of frontier \LLM{}s
or an artefact of specific training choices. No systematic cross-provider
comparison has been conducted.

We address both gaps by applying the identical experimental protocol to four
next-generation models spanning three providers:
\begin{itemize}
  \item \textbf{Claude~Sonnet~4.6} (Anthropic) --- successor to Claude~3.5~Sonnet;
  \item \textbf{Gemini~2.5~Flash} (Google) --- high-throughput frontier model;
  \item \textbf{Gemini~3.1~Pro} (Google) --- Google's premium-tier model;
  \item \textbf{GPT-5.4~Mini} (OpenAI) --- successor to ChatGPT-4o.
\end{itemize}

We evaluate four hypotheses, stated in advance of analysis:
\begin{description}
  \item[\textbf{H1}] \emph{Cooperative bias persists.} Next-generation models
    will maintain or increase the cooperative bias documented by Willis et al.,
    reflecting continued improvements in value alignment.
  \item[\textbf{H2}] \emph{Aggressive capability parity.} Better reasoning will
    enable more effective aggressive strategies, reducing the payoff advantage
    that cooperative agents currently hold over aggressive ones.
  \item[\textbf{H3}] \emph{Cross-provider divergence.} Differences in cooperative
    tendencies across providers will be large enough to attribute to distinct
    training and alignment choices rather than sampling noise.
  \item[\textbf{H4}] \emph{Noise robustness improves.} Newer models will be
    more robust to action noise, particularly Claude, which exhibited marked
    noise sensitivity in the original benchmark.
\end{description}

Three contributions follow from this design. First, we extend the \citet{Willis2025_llm_ipd}
benchmark to four new frontier models and three providers, enabling the first
cross-provider comparative analysis; H1 is confirmed in 10 of 12 model-prompt
combinations, establishing cooperative bias as a broadly shared property of
current frontier \LLM{}s. Second, we introduce the Index of Differential
Capabilities (\ICD), a scalar summary of the aggressive--cooperative payoff gap;
across conditions, \ICD ranges from 0.454 to 0.925, with Gemini~3.1~Pro under
Self-Refine prompting reaching the highest value in the dataset. Third,
we provide a longitudinal comparison within the Anthropic and OpenAI lineages;
provider identity is the strongest predictor of equilibrium outcome---Gemini~2.5
Flash produces aggressive equilibria in up to 77\% of biased-population runs,
while GPT-5.4~Mini reaches 70\% cooperative equilibria under Refine prompting,
confirming H3 and suggesting that alignment choices, not architectural generation
alone, drive variation.

% --- roadmap ---
The remainder of this paper is organised as follows.
Section~\ref{sec:related} reviews related work.
Section~\ref{sec:method} summarises the experimental protocol.
Section~\ref{sec:results} presents our findings.
Section~\ref{sec:discussion} evaluates the four hypotheses and discusses
implications for \MAS design.
Section~\ref{sec:conclusion} concludes.
